\phantomsection
\addcontentsline{toc}{section}{Lampiran B --- Pseudocode Pipeline Migrasi PHP}
\section*{Pseudocode Pipeline Migrasi PHP}
\label{app:pseudocode-pipeline}


Lampiran ini menyajikan pseudocode setiap modul pada pipeline migrasi PHP yang
dikembangkan. Pseudocode merupakan representasi logika program dalam bahasa
semi-formal yang lebih mudah dibaca dibandingkan kode program sesungguhnya,
tanpa terikat pada sintaks bahasa pemrograman tertentu.

\subsection*{Panduan Membaca Pseudocode}

Berikut keterangan notasi yang digunakan pada seluruh pseudocode di lampiran ini.

\begin{center}
\renewcommand{\arraystretch}{1.3}
\begin{tabular}{|p{5.5cm}|p{9cm}|}
\hline
\textbf{Notasi} & \textbf{Keterangan} \\
\hline
\texttt{PROGRAM: NamaProgram} &
  Nama modul/program yang didefinisikan. \\
\hline
\texttt{FUNCTION nama(param):} \newline
\texttt{...END FUNCTION} &
  Definisi fungsi beserta parameternya. Kode di dalam blok ini dijalankan
  saat fungsi dipanggil. \\
\hline
\texttt{param <- nilai} &
  Penugasan nilai ke variabel (dibaca: \textit{param diisi dengan nilai}). \\
\hline
\texttt{IF kondisi THEN} \newline
\texttt{...ELSE...END IF} &
  Percabangan: blok dalam \texttt{THEN} dijalankan jika kondisi benar,
  blok \texttt{ELSE} jika salah. \\
\hline
\texttt{FOR each item in daftar DO} \newline
\texttt{...END FOR} &
  Perulangan: blok dijalankan satu kali untuk setiap elemen dalam daftar. \\
\hline
\texttt{RETURN nilai} &
  Mengembalikan hasil fungsi ke pemanggil. \\
\hline
\texttt{DATA STRUCTURE Nama:} &
  Definisi struktur data (kumpulan atribut yang membentuk satu entitas). \\
\hline
\texttt{CONSTANT NAMA: nilai} &
  Nilai tetap yang tidak berubah selama program berjalan. \\
\hline
\texttt{// komentar} &
  Keterangan tambahan, \textbf{tidak dieksekusi} oleh program. \\
\hline
\texttt{NULL} &
  Tidak ada nilai / kosong. \\
\hline
\texttt{A -> B} &
  Menunjukkan keluaran atau pemetaan: \textit{A menghasilkan B}. \\
\hline
\texttt{A <- B} (di luar penugasan) &
  Menunjukkan asal data: \textit{A berasal dari B}. \\
\hline
\end{tabular}
\end{center}

\bigskip

% --------------------------------------------------------
\section*{B.1 Modul \texttt{main.py} --- Orkestrator Pipeline Utama}
% --------------------------------------------------------

Modul ini adalah titik masuk (\textit{entry point}) program. Ketika pengguna
menjalankan pipeline dari \textit{command line}, modul inilah yang pertama kali
dieksekusi. Tugasnya adalah membaca argumen yang diberikan pengguna (seperti
lokasi folder PHP yang akan dimigrasi, model AI yang dipakai, dan opsi lainnya),
kemudian memanggil fungsi \texttt{run\_pipeline()} yang mengorkestrasi seluruh
tahapan migrasi secara berurutan. Di akhir proses, program keluar dengan kode
\texttt{0} jika sistem dinyatakan patuh ISO, atau kode \texttt{1} jika belum.

\begin{lstlisting}[style=pseudocodestyle]
PROGRAM: MigrationPipeline

FUNCTION main():
    args <- parse_command_line_arguments()
        // argumen: --input, --output, --reports, --skip-ai, --model,
        //          --compare-all, --target-php, --php-version, dll.

    IF output_dir is NOT empty THEN
        prompt user "Hapus isi output/? [Y/n]"
        IF user confirms THEN
            delete contents of output_dir
        ELSE
            exit program
        END IF
    END IF

    result <- run_pipeline(args)

    IF result.iso_status == COMPLIANT THEN
        exit with code 0   // sistem dinyatakan patuh ISO
    ELSE
        exit with code 1   // masih terdapat temuan keamanan
    END IF
END FUNCTION


FUNCTION run_pipeline(input_dir, output_dir, reports_dir, options):
    record start_time

    // Step 1: Pindai keamanan kode SEBELUM konversi (kondisi awal)
    pre_scan <- run_scan(input_dir)

    // Step 2: Konversi sintaks PHP lama ke PHP 8.x menggunakan Rector
    conversion <- run_conversion(input_dir, output_dir, target_version)

    // Step 2b: Analisis dependensi library di composer.json
    composer_analysis <- run_composer_analysis(input_dir, 
    target_version)

    // Step 3: Pindai keamanan kode SESUDAH konversi (kondisi akhir)
    IF conversion succeeded THEN
        post_scan <- run_scan(output_dir)
    END IF

    // Step 4: Analisis statis PHPStan untuk mendeteksi kesalahan tipe
    //         dan referensi kode yang tidak valid
    IF conversion succeeded THEN
        analysis <- run_analysis(output_dir, phpstan_level, 
        php_version)
    END IF

    // Step 5: Analisis AI untuk rekomendasi perbaikan kerentanan
    IF skip_ai == FALSE THEN
        IF compare_all == TRUE THEN
            // Mode perbandingan: uji semua model LLM, 
            3 kali per temuan
            run_llm_comparison(pre_scan.findings, all_models, 
            3_runs_each)
        ELSE
            // Mode normal: gunakan satu model terpilih
            ai_recs <- run_ai_analysis(pre_scan.findings, 
            selected_model)
        END IF
    END IF

    // Step 6: Petakan semua temuan ke kontrol ISO/IEC 27001:2022
    iso_report <- run_iso_mapping(
        pre_scan, analysis, ai_recs, output_path)

    total_duration <- now() - start_time

    print_summary_table(pre_scan, post_scan, conversion, analysis, 
    iso_report)
    save_pipeline_result_json(reports_dir, timestamp)

    RETURN PipelineResult(pre_scan, conversion, post_scan, analysis,
                          ai_recs, iso_report, composer_analysis, 
                          duration)
END FUNCTION
\end{lstlisting}

% --------------------------------------------------------
\section*{B.2 Modul \texttt{scanner.py} --- Pemindaian Keamanan (Semgrep)}
% --------------------------------------------------------

Modul ini bertanggung jawab menjalankan \textit{tool} Semgrep untuk memindai
seluruh file PHP dalam sebuah direktori dan mendeteksi potensi kerentanan
keamanan. Semgrep adalah alat analisis kode statis berbasis pola
(\textit{pattern matching}): ia mencocokkan kode sumber dengan aturan-aturan
yang mendefinisikan pola kode berbahaya. Setiap temuan kemudian diklasifikasikan
ke dalam jenis kerentanan (misalnya SQL Injection atau XSS), diberi tingkat
prioritas dari 1 (kritis) hingga 6 (rendah), dan dipetakan ke kontrol
ISO/IEC 27001:2022 yang relevan.

\begin{lstlisting}[style=pseudocodestyle]
PROGRAM: SemgrepScanner

DATA STRUCTURE ScanFinding:
    // Satu temuan kerentanan dari hasil Semgrep
    rule_id, file_path, line_start, line_end
    message, severity // ERROR | WARNING | INFO
    code_snippet  // potongan kode yang bermasalah
    vuln_type     // SQL Injection | XSS | Deprecated Function | ...
    iso_controls  // kontrol ISO terkait, e.g. ["A.8.28", "A.8.26"]
    priority      // 1..6 (1 = tertinggi / paling kritis)

DATA STRUCTURE ScanResult:
  findings   : list of ScanFinding // diurutkan berdasarkan prioritas
  summary    : ScanSummary         // jumlah total per jenis/severity
  raw_output : JSON dict           // keluaran mentah dari Semgrep

CONSTANT VULNERABILITY_RULES:
    // Pemetaan jenis kerentanan ke prioritas dan kontrol ISO
    priority 1: SQL Injection        -> ISO [A.8.28, A.8.26]
    priority 2: XSS                  -> ISO [A.8.28, A.8.26]
    priority 3: Deprecated Function  -> ISO [A.8.25, A.8.28]
    priority 4: Hardcoded Credentials -> ISO [A.5.17, A.8.28]
    priority 5: SSRF / Path Traversal -> ISO [A.8.28, A.8.29]
    priority 6: Weak Cryptography    -> ISO [A.8.24, A.8.28]


FUNCTION run_scan(target_path):
// Pilih ruleset Semgrep: lokal (reproducible) atau registry online
    IF local rules/ directory exists THEN
        rulesets <- local rules/ directory
    ELSE
        rulesets <- ["p/php", "p/owasp-top-ten"]
    END IF

    cmd <- build_semgrep_command(target_path, rulesets, "--json")
    raw_json <- invoke_subprocess(cmd, timeout=300s)

    findings <- []
    FOR each result in raw_json["results"] DO
        rule_id  <- result.check_id
        severity <- result.extra.severity
        message  <- result.extra.message
        snippet  <- result.extra.lines

      // Klasifikasi ke jenis kerentanan, kontrol ISO, dan prioritas
      (vuln_type, iso_controls, priority) <- classify(rule_id, message)

        findings.append(ScanFinding(...))
    END FOR

    // Urutkan dari yang paling kritis
    findings.sort(by priority ASC, then severity, then file_path)

    summary <- aggregate(findings)
    print_summary_table(summary)

    RETURN ScanResult(findings, summary, raw_json)
END FUNCTION


FUNCTION classify(rule_id, message):
// Cocokkan ID aturan dan pesan ke dalam VULNERABILITY_RULES
  haystack <- lowercase(rule_id + message[:100])

  FOR each (keywords, label, controls, priority) in 
            VULNERABILITY_RULES DO
        IF any keyword in haystack THEN
            RETURN (label, controls, priority)
        END IF
    END FOR

    RETURN ("Other", ["A.8.28"], priority=99)  
    // kategori tidak dikenal
END FUNCTION
\end{lstlisting}

% --------------------------------------------------------
\section*{B.3 Modul \texttt{converter.py} --- Konversi PHP (Rector)}
% --------------------------------------------------------

Modul ini melakukan konversi otomatis kode PHP lama ke PHP 8.x menggunakan
\textit{tool} Rector. Rector adalah \textit{tool} refaktorisasi PHP berbasis
aturan (\textit{rule-based}): ia membaca kode sumber, membangun representasi
pohon sintaks (\textit{Abstract Syntax Tree}), menerapkan aturan transformasi
yang sesuai, lalu menuliskan kode yang sudah diperbarui. Modul ini terlebih
dahulu mendeteksi versi PHP tiap file (PHP 5 atau PHP 7), memilih rantai aturan
migrasi yang tepat, kemudian memverifikasi file mana yang benar-benar berubah
menggunakan sidik jari SHA-256.

\begin{lstlisting}[style=pseudocodestyle]
PROGRAM: RectorConverter

DATA STRUCTURE ConversionFile:
    source_path, output_path
    php_version      // PHP5 | PHP7 | UNKNOWN
    status           // CONVERTED | SKIPPED | FAILED
    changes_count    // jumlah perubahan yang diterapkan
    applied_rectors  // daftar nama aturan Rector yang aktif
    iso_controls     // ["A.8.25"] - Secure Development Lifecycle

DATA STRUCTURE ConversionResult:
    files             : list of ConversionFile
    summary           : ConversionSummary
    rector_raw_output : string (JSON)  // keluaran mentah Rector

CONSTANT LEVEL_SETS_PHP5:
    // Rantai migrasi penuh untuk kode PHP 5 (setiap versi minor)
    UP_TO_PHP_53, UP_TO_PHP_54, ..., UP_TO_PHP_85

CONSTANT LEVEL_SETS_PHP7:
    // Rantai lebih pendek untuk kode PHP 7
    UP_TO_PHP_70, ..., UP_TO_PHP_85

CONSTANT QUALITY_SETS:
    // Aturan peningkatan kualitas kode (selalu diterapkan)
    CODE_QUALITY, DEAD_CODE, EARLY_RETURN, TYPE_DECLARATION


FUNCTION run_conversion(input_path, output_path, target_version):
    php_files <- collect_all_php_files(input_path)

// Step 1: Deteksi versi PHP per file dari pola kode yang digunakan
    version_map <- {}
    FOR each file in php_files DO
        version_map[file] <- detect_php_version(file)
    END FOR

// Step 2: Pilih rantai aturan Rector sesuai versi PHP yang 
    ditemukan
    IF PHP5 in version_map.values() THEN
        level_sets <- LEVEL_SETS_PHP5 truncated at target_version
    ELSE
        level_sets <- LEVEL_SETS_PHP7 truncated at target_version
    END IF
    all_sets <- level_sets + QUALITY_SETS

// Step 3: Salin input/ ke output/ agar input asli tidak pernah 
    diubah
    copy_tree(input_path, output_path)

// Step 4: Catat hash SHA-256 semua file SEBELUM Rector dijalankan
    pre_hashes <- snapshot_sha256(output_path)

// Step 5: Tulis konfigurasi Rector sementara ke direktori temp 
    sistem
    config_path <- write_rector_config(output_path, all_sets)

// Step 6: Jalankan Rector sebagai subprocess
    (exit_code, stdout, stderr) <- invoke_rector(config_path, 
    timeout=600s)

// Step 7: Catat hash SHA-256 semua file SESUDAH Rector
    post_hashes <- snapshot_sha256(output_path)

// Step 8: Hapus konfigurasi sementara
    delete(config_path)

// Step 9: Tentukan status tiap file dengan membandingkan hash
    rector_data <- parse_rector_json(stdout)
    FOR each source_file in php_files DO
        pre  <- pre_hashes[output_copy_of(source_file)]
        post <- post_hashes[output_copy_of(source_file)]

        IF file is in rector error list THEN
            status <- FAILED     // Rector gagal memproses file ini
        ELSE IF pre != post THEN
            status <- CONVERTED  // file berhasil dikonversi
        ELSE
            status <- SKIPPED    // tidak ada perubahan yang perlu 
            diterapkan
        END IF

        files.append(ConversionFile(source_file, status, ...))
    END FOR

    summary <- aggregate(files, elapsed_time)
    print_summary_table(summary)

    RETURN ConversionResult(files, summary, stdout)
END FUNCTION


FUNCTION detect_php_version(file_path):
    // Baca 8KB pertama file untuk efisiensi
    content <- read first 8KB of file

    // Cek pola fungsi/sintaks yang hanya ada di PHP 5
    IF any PHP5 pattern matches (mysql_connect, ereg, split) THEN
        RETURN PHP5
    END IF

    // Cek pola yang diperkenalkan di PHP 7
    IF any PHP7 pattern matches (fn =>, ??, intdiv) THEN
        RETURN PHP7
    END IF

    RETURN UNKNOWN
END FUNCTION
\end{lstlisting}

% --------------------------------------------------------
\section*{B.4 Modul \texttt{analyzer.py} --- Analisis Statis (PHPStan)}
% --------------------------------------------------------

Modul ini menjalankan PHPStan, sebuah alat analisis statis untuk PHP yang
memeriksa kode tanpa mengeksekusinya. PHPStan mendeteksi kesalahan seperti
pemanggilan fungsi yang tidak ada, ketidakcocokan tipe data, kode yang tidak
pernah dieksekusi (\textit{dead code}), dan referensi variabel yang tidak valid.
Karena proyek target menggunakan kerangka kerja CodeIgniter 3 (CI3) yang
mengandalkan konvensi \textit{magic method} dan konstanta global yang tidak
dikenali PHPStan, modul ini secara khusus menyaring pesan \textit{false positive}
dari CI3 sebelum hasilnya dilaporkan.

\begin{lstlisting}[style=pseudocodestyle]
PROGRAM: PHPStanAnalyzer

DATA STRUCTURE AnalysisError:
    file_path, line
    message, severity   // ERROR | WARNING
    iso_controls        // kontrol ISO yang relevan

DATA STRUCTURE AnalysisResult:
    errors     : list of AnalysisError  // diurutkan per file lalu 
                 per baris
    summary    : AnalysisSummary
    raw_output : JSON dict

CONSTANT CI3_FRAMEWORK_NOISE:
    // False positive khas arsitektur CodeIgniter 3 -- diabaikan:
    "extends unknown class CI_*"
    "Constant ENVIRONMENT/FCPATH/APPPATH not found"
    "$this->db, $this->input, $this->session, ..."
    "Function base_url/redirect/form_open not found"


FUNCTION run_analysis(target_path, level, php_version):
    // Cari binary PHPStan (vendor/bin/ atau PATH global)
    phpstan_exe <- find_phpstan(vendor/bin/ OR global PATH)

    cmd <- build_phpstan_command(target_path, level, php_version,
                            "--error-format=json", "--no-progress")
    (exit_code, stdout, stderr) <- invoke_subprocess(cmd, timeout=300s)

    raw_json <- parse_phpstan_json(stdout)

    errors <- []
    noise_count <- 0
    FOR each (file_path, file_data) in raw_json["files"] DO
        FOR each msg_entry in file_data["messages"] DO
            message <- msg_entry.message
            line    <- msg_entry.line

            // Lewati pesan yang merupakan false positive CI3
            IF message matches CI3_FRAMEWORK_NOISE THEN
                noise_count <- noise_count + 1
                CONTINUE
            END IF

            (severity, iso_controls) <- classify_error(message)
            errors.append(AnalysisError(file_path, line, message,
                                        severity, iso_controls))
        END FOR
    END FOR

    errors.sort(by file_path, then line)

    summary <- aggregate(errors, noise_count, elapsed_time)
    print_results_table(errors, summary)

    RETURN AnalysisResult(errors, summary, raw_json)
END FUNCTION


FUNCTION classify_error(message):
    lower_msg <- lowercase(message)

    // Kode mati atau jalur eksekusi yang tidak dapat dicapai
    IF "dead code" OR "unreachable" OR "never returns" 
        in lower_msg THEN
        RETURN ("WARNING", ["A.8.25"])
    END IF

    // Referensi ke fungsi, kelas, atau konstanta yang tidak ada
    IF "undefined" OR "not found" OR "unknown class" in lower_msg THEN
        RETURN ("ERROR", ["A.8.28", "A.8.25"])
    END IF

    // Ketidakcocokan tipe data atau nilai null yang tidak ditangani
  IF "type" OR "null" OR "incompatible" OR "expects" in lower_msg THEN
        RETURN ("ERROR", ["A.8.28"])
    END IF

    RETURN ("ERROR", ["A.8.28"])   
    // fallback untuk pesan yang tidak dikenal
END FUNCTION
\end{lstlisting}

% --------------------------------------------------------
\section*{B.5 Modul \texttt{ai\_engine.py} --- Analisis AI (LLM via Ollama)}
% --------------------------------------------------------

Modul ini mengintegrasikan model bahasa besar (\textit{Large Language Model} /
LLM) lokal yang berjalan melalui Ollama untuk memberikan rekomendasi perbaikan
kode terhadap setiap kerentanan yang ditemukan Semgrep. Modul mendukung dua
mode: mode normal (satu model) dan mode perbandingan (lima model, tiga kali
pengujian per temuan). Setiap potongan kode rentan dikirim ke LLM beserta
konteks jenis kerentanan dan kontrol ISO yang berlaku, kemudian respons LLM
diurai untuk mengekstrak penjelasan, kode perbaikan, dan skor kepercayaan
(\textit{confidence score}).

\begin{lstlisting}[style=pseudocodestyle]
PROGRAM: PHPAIEngine

DATA STRUCTURE AIRecommendation:
    original_code, suggested_fix, explanation
    confidence    // skor kepercayaan LLM: 0.0 - 1.0
    iso_controls, vuln_type
    file_path, line_start
    model_used, fim_used  // apakah menggunakan Fill-in-Middle

CONSTANT COMPARISON_MODELS:
    ["deepseek-coder:6.7b", "qwen2.5-coder:7b",
     "codellama:7b", "mistral:7b", "llama3.1:8b"]

// Hanya proses temuan dengan prioritas 1-5 (abaikan kategori "Other")
CONSTANT PRIORITY_THRESHOLD: 5
// Batasi panjang kode yang dikirim ke LLM (karakter)
CONSTANT MAX_CODE_CHARS: 2000


FUNCTION run_ai_analysis(findings, model):
    // Saring hanya temuan yang cukup kritis untuk dianalisis AI
    eligible <- [f for f in findings WHERE f.priority <= 
    PRIORITY_THRESHOLD]

    // Pastikan server Ollama lokal sedang berjalan
    IF Ollama NOT reachable at localhost:11434 THEN
        print warning "Ollama offline -- tahapan AI dilewati"
        RETURN []
    END IF

    recommendations <- []
    FOR each finding in eligible DO
        rec <- analyze_snippet(finding.code_snippet, finding.vuln_type,
                               finding.iso_controls, finding.file_path,
                               finding.line_start, model)
        recommendations.append(rec)
    END FOR

    print_batch_summary(recommendations)
    RETURN recommendations
END FUNCTION


FUNCTION analyze_snippet(code, vuln_type, iso_controls,
                         file_path, line_start, model):
    truncated <- code[:MAX_CODE_CHARS]

    // Kode pendek (<= 3 baris): gunakan Fill-in-Middle (FIM) prompt
    // Kode lebih panjang: gunakan chat biasa
    IF lines(truncated) <= 3 THEN
      prompt   <- build_fim_prompt(truncated, vuln_type, iso_controls)
      fim_used <- TRUE
    ELSE
    messages <- build_chat_messages(truncated, vuln_type, iso_controls)
    fim_used <- FALSE
    END IF

    // temperature=0.1: output deterministik, minim variasi acak
    raw_text <- call_ollama_chat(messages, model, temperature=0.1)

    RETURN parse_response(raw_text, truncated, iso_controls,
                          vuln_type, file_path, line_start, fim_used)
END FUNCTION


FUNCTION build_chat_messages(code, vuln_type, iso_controls):
    // Kondisi A (standar): prompt identik untuk semua model
   system_msg <- "You are a PHP security and migration expert
                   controls [{iso_controls}]."
    user_msg   <- "TASK: Analyze [{vuln_type}] vulnerability...
                   VULNERABLE CODE: (kode PHP disisipkan di sini)
                   1. Explain the vulnerability.
                   2. Provide corrected PHP 8.x code.
                   Last line: CONFIDENCE: [0-100]"
    RETURN [system_msg, user_msg]
END FUNCTION


FUNCTION build_optimized_chat_messages(code, vuln_type, iso_controls, 
model):
    // Kondisi B: prompt dioptimalkan khusus per model
    IF model == "qwen2.5-coder:7b" THEN
        // format diperlunak; CONFIDENCE tidak diwajibkan
    ELSE IF model == "deepseek-coder:6.7b" THEN
        // instruksi pendek dan langsung
    ELSE IF model == "mistral:7b" THEN
        // satu contoh few-shot untuk mengajarkan format keluaran
    ELSE IF model == "llama3.1:8b" THEN
        // role-play persona kuat + 4 langkah prosedural
    ELSE IF model == "codellama:7b" THEN
        // sama dengan Kondisi A + instruksi validitas sintaks php -l
    ELSE
        // fallback ke Kondisi A standar
    END IF
    RETURN [system_msg, user_msg]
END FUNCTION


FUNCTION parse_response(raw_text, original_code, iso_controls, ...):
    IF raw_text is empty THEN
        RETURN placeholder_recommendation(confidence=0.0)
    END IF

    // --- Ekstraksi penjelasan ---
    // Coba tag XML <explanation>, lalu "1. EXPLANATION:", lalu 
       teks bebas

    // --- Ekstraksi kode perbaikan ---
    // Coba tag XML <fix>, 
    // lalu "2. FIXED CODE:", lalu blok ```php ... ```

    // --- Ekstraksi skor kepercayaan (confidence) ---
    // Coba tag <confidence>N</confidence>
    // Coba kata kunci "CONFIDENCE: N" (integer 0-100)
    // Coba desimal 0.x di dekat kata "confidence"
    // Coba persentase di dekat kata "confidence"
    // Coba bahasa natural: "not sure"->0.30, "confident"->0.75
    // Jika tidak ditemukan -> kembalikan 0.0

    RETURN AIRecommendation(original_code, suggested_fix, explanation,
                            confidence, iso_controls, ...)
END FUNCTION


FUNCTION run_llm_comparison(findings, models, runs_per_finding,
                             prompt_mode, reports_dir):
    // Saring temuan yang memenuhi ambang prioritas
    eligible <- [f for f in findings WHERE 
                f.priority <= PRIORITY_THRESHOLD]

    report <- {models_compared, runs_per_finding, prompt_mode, 
               findings: []}

    FOR each model in models DO
        IF Ollama NOT reachable THEN
            mark model as unavailable
            CONTINUE
        END IF

        FOR each finding in eligible DO
            IF prompt_mode == "standard" THEN
             // Kondisi A: prompt identik untuk semua model
                messages <- build_chat_messages(finding)
            ELSE
               // Kondisi B: prompt dioptimalkan per model
               messages <- build_optimized_chat_messages
                           (finding, model)
            END IF

            confidences <- []
            times       <- []
            FOR run_index = 1 to runs_per_finding DO
                (raw_text, elapsed) <- call_ollama_chat_timed
                                       (messages, model)
                confidence <- extract_confidence(raw_text)
                format_ok  <- check_format_valid(raw_text)
                confidences.append(confidence)
                times.append(elapsed)
            END FOR

            // Statistik agregat per temuan per model
            finding_result <- {
            confidence_mean       : mean(confidences),
            confidence_variance   : variance(confidences),
            confidence_stdev      : stdev(confidences),
            format_compliance_rate: count(format_ok)/runs_per_finding,
            mean_inference_time_sec: mean(times)
            }
            report.append(finding_result)
        END FOR
    END FOR

    save_json(report, reports_dir / "llm_comparison_{timestamp}.json")
    print_comparison_table(report)
    RETURN report
END FUNCTION
\end{lstlisting}

% --------------------------------------------------------
\section*{B.6 Modul \texttt{iso\_mapper.py} --- Pemetaan ISO/IEC 27001:2022}
% --------------------------------------------------------

Modul ini adalah "hakim akhir" pipeline yang mengumpulkan seluruh temuan dari
Semgrep, PHPStan, dan rekomendasi AI, kemudian memetakannya ke enam kontrol
keamanan ISO/IEC 27001:2022 Annex A yang relevan. Setiap kontrol diberi status:
\textit{COMPLIANT} (tidak ada temuan), \textit{PARTIAL} (hanya peringatan), atau
\textit{NON\_COMPLIANT} (terdapat kesalahan kritis). Status terburuk di antara
semua kontrol menjadi status keseluruhan sistem. Rekomendasi AI dicatat sebagai
bukti (\textit{evidence}) namun \textbf{tidak} memengaruhi penentuan status
kepatuhan.

\begin{lstlisting}[style=pseudocodestyle]
PROGRAM: ISOMapper

DATA STRUCTURE ISOControl:
    control_id    // e.g. "A.8.28"
    title, description
    status        // COMPLIANT | PARTIAL | NON_COMPLIANT
    findings_count
    evidence      // daftar string bukti temuan (maks. 30 entri)

DATA STRUCTURE ISOReport:
    controls          : dict of ISOControl (kunci: control_id)
    overall_status    : status terburuk dari semua kontrol
    total_findings    : total temuan (Semgrep + PHPStan)
    critical_findings : jumlah temuan ERROR-severity
    generated_at      : waktu laporan dibuat

CONSTANT CONTROL_REGISTRY:
    A.5.17  Authentication Information (pengelolaan kredensial)
    A.8.24  Use of Cryptography (penggunaan enkripsi)
    A.8.25  Secure Development Lifecycle (siklus pengembangan aman)
    A.8.26  Application Security Requirements 
            (persyaratan keamanan aplikasi)
    A.8.28  Secure Coding (praktik penulisan kode aman)
    A.8.29  Security Testing in Development  (pengujian keamanan)

STATUS RULES:
    NON_COMPLIANT : terdapat >= 1 temuan ERROR pada kontrol ini
    PARTIAL       : hanya temuan WARNING, tidak ada ERROR
    COMPLIANT     : tidak ada temuan sama sekali pada kontrol ini


FUNCTION run_iso_mapping(scan_result, analysis_result, ai_recs, 
                         output_path):
    // Step 1: Inisialisasi wadah temuan untuk setiap kontrol
    findings_map <- {}
    FOR each control in CONTROL_REGISTRY DO
        findings_map[control] <- []
    END FOR

    // Step 2: Masukkan temuan Semgrep ke peta kontrol
    IF scan_result != NULL THEN
        FOR each finding in scan_result.findings DO
            // INFO dinaikkan ke WARNING agar tidak ada status 
            // COMPLIANT palsu
            severity <- normalize(finding.severity)
            evidence <- "[Semgrep][vuln_type] file:baris -- pesan"
            FOR each ctrl in finding.iso_controls DO
                findings_map[ctrl].append((severity, evidence))
            END FOR
        END FOR
    END IF

    // Step 3: Masukkan kesalahan PHPStan ke peta kontrol
    IF analysis_result != NULL THEN
        FOR each error in analysis_result.errors DO
            IF error.file_path == "<phpstan>" THEN CONTINUE END IF
            evidence <- "[PHPStan] file:baris -- pesan"
            FOR each ctrl in error.iso_controls DO
                findings_map[ctrl].append((severity, evidence))
            END FOR
        END FOR
    END IF

    // Step 4: Kumpulkan bukti rekomendasi AI (hanya informatif)
    ai_evidence_map <- {}
    FOR each rec in ai_recs DO
        evidence <- "[AI/model][vuln_type] file:baris -- fix 
                     (confidence)"
        FOR each ctrl in rec.iso_controls DO
            ai_evidence_map[ctrl].append(evidence)
        END FOR
    END FOR

    // Step 5: Bangun objek ISOControl untuk setiap kontrol
    controls <- {}
    FOR each (ctrl_id, title, description) in CONTROL_REGISTRY DO
        entries <- findings_map[ctrl_id]
        ai_evs  <- ai_evidence_map[ctrl_id]
        status  <- determine_status(entries)
        // Gabungkan bukti ERROR + WARNING + AI, batasi 30 entri
        evidence <- (ERROR entries + WARNING entries + ai_evs)[:30]
        controls[ctrl_id] <- ISOControl(ctrl_id, title, description,
                                status, count(entries), evidence)
    END FOR

    // Step 6: Status keseluruhan = status terburuk di 
    // antara semua kontrol
    overall_status <- max(ctrl.status for ctrl in controls)

    report <- ISOReport(controls, overall_status,
                        total_findings, critical_findings, now())

    print_control_table(report)
    save_json(report, output_path)
    RETURN report
END FUNCTION


FUNCTION determine_status(entries):
    severities <- {sev for (sev, _) in entries}
    IF "ERROR"   in severities THEN RETURN NON_COMPLIANT END IF
    IF "WARNING" in severities THEN RETURN PARTIAL       END IF
    RETURN COMPLIANT
END FUNCTION
\end{lstlisting}

% --------------------------------------------------------
\section*{B.7 Modul \texttt{composer\_analyzer.py} --- Analisis Dependensi}
% --------------------------------------------------------

Composer adalah manajer dependensi standar PHP. File \texttt{composer.json}
mendaftarkan semua \textit{library} pihak ketiga yang digunakan proyek.
Modul ini membaca file tersebut dan mengevaluasi kompatibilitas setiap
\textit{library} terhadap PHP 8.x menggunakan basis data paket yang sudah
diketahui. Hasilnya berupa rekomendasi: apakah \textit{library} aman digunakan
(\textit{safe}), perlu diperbarui versinya (\textit{upgrade}), perlu diganti
dengan alternatif yang didukung (\textit{replace}), atau perlu diperiksa
secara manual karena datanya belum tersedia (\textit{manual\_review}).

\begin{lstlisting}[style=pseudocodestyle]
PROGRAM: ComposerAnalyzer

DATA STRUCTURE DependencyRecommendation:
    package                // nama library, e.g. "guzzlehttp/guzzle"
    current_version_constraint  // versi yang digunakan, e.g. "^6.0"
    status    // "safe" | "upgrade" | "replace" | "manual_review"
    suggestion  // teks rekomendasi aksi yang harus dilakukan
    iso_control // "A.8.25" - Secure Development Lifecycle

DATA STRUCTURE ComposerAnalysisResult:
    composer_found    : bool
    php_constraint    // e.g. ">=7.4" dari composer.json
    recommendations   : list of DependencyRecommendation
    issues_count      // jumlah library yang berstatus bukan "safe"

CONSTANT KNOWN_PACKAGES:
    // Basis data kompatibilitas library terhadap PHP 8.x
    phpoffice/phpexcel      -> replace  ke phpoffice/phpspreadsheet
    dompdf/dompdf           -> upgrade  ke versi >=2.0
    swiftmailer/swiftmailer -> replace  ke symfony/mailer
    codeigniter/framework   -> replace  ke codeigniter4/framework
    guzzlehttp/guzzle       -> upgrade  ke versi >=7.0
    monolog/monolog         -> safe     (>=2.0 mendukung PHP 8.x)
    firebase/php-jwt        -> safe     (>=6.0 mendukung PHP 8.x)
    ... (dan lainnya)


FUNCTION run_composer_analysis(input_dir, target_php):
    composer_path <- input_dir / "composer.json"

    // Jika proyek tidak menggunakan Composer, lewati tahapan ini
    IF composer_path does NOT exist THEN
        RETURN ComposerAnalysisResult(composer_found=FALSE)
    END IF

    data <- parse_json(composer_path)
    php_constraint <- data["require"]["php"]  // e.g. ">=7.4 <8.0"

    // Kumpulkan semua dependensi (runtime + development),
    // kecuali entri meta "php" dan "ext-*"
    deps <- merge(data["require"], data["require-dev"])
    deps <- filter(deps, exclude php and ext-*)

    recommendations <- []
    FOR each (package, version) in deps DO
        rec <- evaluate_package(package, version, target_php)
        recommendations.append(rec)
    END FOR

    print_dependency_table(recommendations)
    RETURN ComposerAnalysisResult(composer_found=TRUE,
                                  php_constraint, recommendations)
END FUNCTION


FUNCTION evaluate_package(package, version, target_php):
    IF package in KNOWN_PACKAGES THEN
        (action, suggestion, php8_safe) <- KNOWN_PACKAGES[package]
        status <- "safe" IF php8_safe ELSE action
    ELSE
    // Library tidak ada dalam basis data: minta pemeriksaan manual
        status     <- "manual_review"
        suggestion <- "Tidak ada data PHP 8.x untuk paket ini. "
                    + "Periksa secara manual."
    END IF

RETURN DependencyRecommendation(package, version, status, suggestion)
END FUNCTION
\end{lstlisting}

% --------------------------------------------------------
\section*{B.8 Ringkasan Alur Pipeline}
% --------------------------------------------------------

Berikut adalah gambaran keseluruhan alur eksekusi pipeline dari awal hingga
akhir. Setiap langkah menghasilkan sebuah objek hasil (\textit{result object})
yang digunakan oleh langkah berikutnya. Kode keluar (\textit{exit code})
\texttt{0} berarti sistem dinyatakan patuh ISO, sedangkan \texttt{1} berarti
masih terdapat temuan yang perlu ditangani.

\begin{lstlisting}[style=pseudocodestyle]
INPUT: File PHP sumber (PHP 5.x / 7.x) di direktori input/

Tahap 1  run_scan(input/)
         -> ScanResult  [Semgrep, kondisi SEBELUM konversi]

Tahap 2  run_conversion(input/, output/)
         -> ConversionResult  [Rector: PHP lama ke PHP 8.x]

Tahap 2b run_composer_analysis(input/)
         -> ComposerAnalysisResult  [kompatibilitas library]

Tahap 3  run_scan(output/)
         -> ScanResult  [Semgrep, kondisi SESUDAH konversi]

Tahap 4  run_analysis(output/)
         -> AnalysisResult  [PHPStan: analisis tipe & referensi]

Tahap 5  run_ai_analysis(pre_scan.findings)
         -> list[AIRecommendation]  [Ollama LLM: rekomendasi perbaikan]
   ATAU  run_llm_comparison(pre_scan.findings, 5_models, 3_runs)
         -> LLMComparisonReport  [perbandingan antar model]

Tahap 6  run_iso_mapping(pre_scan, analysis, ai_recs)
         -> ISOReport  [pemetaan ke 6 kontrol ISO/IEC 27001:2022]

OUTPUT:
  output/   -- file PHP 8.x hasil konversi
  reports/
    iso_report_{timestamp}.json          // laporan kepatuhan ISO
    pipeline_result_{timestamp}.json     // ringkasan seluruh pipeline
    llm_comparison_{timestamp}.json      // (jika mode --compare-all)

KODE KELUAR:
  0 -> ISO overall_status == COMPLIANT   (semua kontrol terpenuhi)
  1 -> NON_COMPLIANT / PARTIAL / ada tahapan yang gagal
\end{lstlisting}
